% Editar este fichero y rellenar con el resumen del TFM
La predicción de glucosa en sangre es una herramienta clave en la gestión de la diabetes mellitus tipo 1 (DM1), ya que permite anticipar episodios de hipoglucemia e hiperglucemia y mejorar la calidad de vida de los pacientes. Los modelos de aprendizaje profundo basados en redes LSTM (\textit{Long Short-Term Memory}) han demostrado un rendimiento sólido en este dominio gracias a su capacidad de capturar dependencias temporales complejas en señales de monitorización continua de glucosa (CGM). Sin embargo, la heterogeneidad demográfica de las cohortes disponibles —en particular las diferencias de edad y sexo— puede introducir sesgos sistemáticos en el entrenamiento, comprometiendo la equidad y la generalización de estos modelos hacia subgrupos infrarrepresentados.

Este Trabajo Fin de Máster parte de la hipótesis de que equilibrar la composición demográfica del conjunto de entrenamiento mejora el rendimiento clínico de un modelo LSTM de predicción de glucosa, especialmente para los subgrupos minoritarios. Para contrastarla, se diseña e implementa un \textit{pipeline} experimental de balanceo aplicado exclusivamente a los pliegues de entrenamiento en un esquema de validación cruzada por grupos (\textit{Group K-Fold}), que garantiza la separación estricta de pacientes entre particiones. Se evalúa una taxonomía de técnicas de balanceo —organizadas en familias aleatorias, guiadas e híbridas— aplicadas sobre dos dimensiones demográficas (edad y sexo) más una condición de control sin balanceo, en tres conjuntos de datos de CGM con características estructurales distintas: \textbf{T1DiabetesGranada} (643 pacientes, más de 30 millones de mediciones), \textbf{REPLACE-BG} y \textbf{DiaTrend}.

El análisis estadístico combina pruebas de Friedman, comparaciones \textit{post-hoc} de Nemenyi, tamaños del efecto de Cohen's~$d$ y agregación de \textit{rankings} mediante el método de Borda, con el fin de determinar si el balanceo demográfico produce efectos significativos y consistentes, y en qué condiciones resulta beneficioso desde una perspectiva clínica.
